\documentclass[11pt,a4paper]{article}
\usepackage[T1]{fontenc}
\usepackage{lmodern}
\usepackage[margin=1in]{geometry}
\usepackage{amsmath,amssymb,amsthm}
\usepackage{microtype}
\usepackage[colorlinks=true,linkcolor=blue,citecolor=blue,urlcolor=blue]{hyperref}
\newtheorem{theorem}{Theorem}
\newtheorem{lemma}[theorem]{Lemma}
\newtheorem{proposition}[theorem]{Proposition}
\newcommand{\F}{\mathbb F_2}
\newcommand{\A}{\mathcal A}
\newcommand{\PH}{\operatorname{PH}}
\newcommand{\ind}{\mathbf 1}
\title{The homology of a persistence fiber without finite bars}
\author{Kaiwen Jin\\[2pt]\small\texttt{yui.ui.siki@gmail.com}}
\date{14 October 2026}

\begin{document}
\maketitle
\begin{abstract}
For a finite connected simplicial complex, consider all simplex-wise filters
whose ordinary persistent homology has one infinite degree-zero bar born at
time~1 and all positive-degree bars born at time~2, with no finite bars.
We prove that this fiber has the same homology over $\F$ as the original
complex. The fiber has a deformation retract homeomorphic to the order
complex of nonempty acyclic subcomplexes of the original complex.
Comparing that poset with the face poset gives the homology
isomorphism and proves the corresponding Betti-number conjecture of
Leygonie and Henselman-Petrusek in the connected setting.
\end{abstract}

\section{The fiber}

A persistence barcode records the homology of a filtration and the maps
between its stages. Different filters can have the same barcode, so the
topology of a barcode fiber describes ambiguity left by this summary.
Leygonie and Henselman-Petrusek \cite{LH} compute such fibers on cell
complexes. Their Conjecture~1, in the preceding connected setting, predicts
that a particular fiber without finite bars has the Betti numbers of the
underlying simplicial complex. Polyhedral descriptions and dimension bounds
for multiparameter fibers are given in \cite[Theorem~5.3 and
Proposition~5.9]{HTT}. Here the absence of finite bars permits a direct
homological comparison.

Let $K$ be a finite connected simplicial complex. Throughout, homology has
coefficients in $\F$ and is ordinary and unreduced unless a tilde is present.
We identify a simplicial complex with its geometric realization when writing
its homology. Simplicial and singular homology agree, including for pairs
\cite[Theorem~2.27 and p.~153]{Hatcher}.
Write $\beta_q(K)=\dim_{\F}H_q(K;\F)$.
Let $D_K$ contain one infinite bar born at~1 in degree zero and
$\beta_q(K)$ infinite bars born at~2 in degree $q\geq1$, with no finite bars.
If all positive-degree Betti numbers vanish set $m=1$, and otherwise set
$m=2$. Put $I=[1,m+1]$.

A filter assigns a value $f(\sigma)\in I$ to every nonempty simplex
$\sigma\in K$, with $f(\tau)\leq f(\sigma)$ whenever $\tau$ is a face of
$\sigma$. Its sublevel complex is
\[
 K_t(f)=\{\sigma\in K:f(\sigma)\leq t\}.
\]
Let $\PH(f)$ be the graded barcode of this filtration, and give
\[
 F_K=\{f:\PH(f)=D_K\}
\]
the topology inherited from $I^{|K|}$, where $|K|$ here denotes the number of
nonempty simplices.

\begin{theorem}\label{thm:main}
For every finite connected simplicial complex $K$ and every integer $q\geq0$,
\[
 H_q(F_K;\F)\cong H_q(K;\F).
\]
In particular, $\dim_{\F}H_q(F_K;\F)=\beta_q(K)$.
\end{theorem}

The filters in this theorem assign values to all simplices. The proof uses
their whole sublevel complexes. Its homological comparison uses acyclic
subcomplexes, which need not be contractible; it does not assert a homotopy
equivalence between $F_K$ and $K$.

Assume for now that $K$ is nonempty. Call a subcomplex \emph{acyclic} if it
is nonempty and its reduced homology over $\F$ vanishes in every degree.
The barcode $D_K$ implies that $K_t(f)$ is acyclic for $1\leq t<2$ and that
every inclusion $K_t(f)\hookrightarrow K$ induces homology isomorphisms for
$t\geq2$. In degree zero these isomorphisms agree with augmentation: the
class of every vertex maps to $1\in\F$.

\begin{lemma}\label{lem:truncate}
The subspace
\[
 G_K=\{f\in F_K:f(\sigma)\leq2\text{ for all }\sigma\in K\}
\]
is a strong deformation retract of $F_K$.
\end{lemma}

\begin{proof}
For $0\leq s\leq1$ define
\[
 h_s(x)=
 \begin{cases}
 x,&x\leq2,\\
 2+(1-s)(x-2),&x\geq2,
 \end{cases}
 \qquad R_s(f)=h_s\circ f.
\]
This is jointly continuous and nondecreasing in $x$, so $R_s(f)$ is a
filter. It fixes all values at most~2. For $s<1$, the function $h_s$ is
strictly increasing and fixes~1 and~2. It only reindexes stages after~2,
whose homology maps to $K$ are already isomorphisms; after its largest value,
the sublevel is $K$. Thus it preserves the entire persistence module up to
isomorphism. At $s=1$, the sublevels below~2 are unchanged and the sublevel
at and after~2 is $K$, so the same barcode is obtained. Finally $R_0$ is the
identity, $R_1(f)=\min(f,2)$, and every $R_s$ fixes $G_K$. If $m=1$, all
values are already at most~2 and this deformation is the identity.
\end{proof}

A filter $g:K\to[1,2]$ belongs to $G_K$ precisely when
\begin{equation}\label{eq:criterion}
 \min_{\sigma\in K}g(\sigma)=1,
 \qquad K_t(g)\text{ is acyclic for every }1\leq t<2.
\end{equation}
Indeed, necessity follows from the barcode. Conversely, the degree-zero
module is constant and one-dimensional from~1 onward, since inclusions
between connected nonempty complexes preserve its augmentation generator.
Positive-degree homology is zero before~2 and equals that of $K$ from~2
onward, because $K_2(g)=K$. This gives exactly $D_K$, also when $m=1$.

\section{Chains of acyclic subcomplexes}

Let $\A(K)$ be the finite poset of nonempty acyclic subcomplexes of $K$,
ordered by inclusion. For a finite poset $P$, write $\Delta P$ for the
simplicial complex whose nonempty simplices are the nonempty chains in $P$.
We write $|\Delta P|$ for its geometric realization.

\begin{proposition}\label{prop:chains}
There is a homeomorphism $|\Delta\A(K)|\to G_K$, affine on every simplex,
given in barycentric coordinates by
\begin{equation}\label{eq:phi}
 \Phi(\lambda)(\sigma)
 =2-\sum_{A\in\A(K)}\lambda_A\ind_{\{\sigma\in A\}}.
\end{equation}
Here $\lambda_A\geq0$, $\sum_A\lambda_A=1$, and the positive support is a
chain of subcomplexes.
\end{proposition}

\begin{proof}
The formula is continuous and gives values in $[1,2]$. If $\tau$ is a face
of $\sigma$, every subcomplex containing $\sigma$ contains $\tau$, so the
value on $\tau$ is at most the value on $\sigma$.

Write the positive support as $A_1\subsetneq\cdots\subsetneq A_r$, with
coefficients $\lambda_i>0$. Every simplex in $A_1$ has value~1. A simplex
whose first membership is in $A_j$ has value
$2-\sum_{i=j}^r\lambda_i$, and a simplex outside $A_r$ has value~2.
The distinct sublevels below~2 are therefore exactly $A_1,\ldots,A_r$.
They are acyclic, so \eqref{eq:criterion} puts $\Phi(\lambda)$ in $G_K$.

For the inverse, list the distinct values of $g\in G_K$ less than~2 as
$1=a_1<\cdots<a_r<2$, and put $a_{r+1}=2$. Define
\begin{equation}\label{eq:inverse}
 A_i=K_{a_i}(g),\qquad \lambda_i=a_{i+1}-a_i
 \quad(1\leq i\leq r).
\end{equation}
These are nonempty acyclic subcomplexes forming a strict chain, and their
positive coefficients sum to~1. If $g(\sigma)=a_j$, its memberships are
exactly those with $i\geq j$, and
\[
 2-\sum_{i=j}^r(a_{i+1}-a_i)=a_j.
\]
A simplex of value~2 belongs to none of these sublevels. Hence
\eqref{eq:inverse} recovers $g$ through \eqref{eq:phi}, and uniquely recovers
the positive support and its coefficients. In particular, zero coefficients
on faces of the order complex create no ambiguity. If $\max g<2$, then
$A_r=K$, which is allowed when $K$ is acyclic; if $\max g=2$, its final
sublevel at~2 is omitted from the list.

Thus $\Phi$ is a continuous bijection from a compact finite simplicial
realization to the Hausdorff subspace $G_K$ of Euclidean space, and is a
homeomorphism \cite[Theorem~26.6]{Munkres}.
\end{proof}

\section{The homology comparison}

The needed comparison is the acyclic case of homological poset fiber
theorems; see \cite[author manuscript, Section~4]{BWW}. We give its finite-poset proof to make
explicit why acyclicity suffices.

\begin{lemma}\label{lem:poset}
Let $u:S\to P$ be an order-preserving map of finite nonempty posets. If
\[
 \Delta\{x\in S:u(x)\leq p\}
\]
is nonempty and acyclic over $\F$ for every $p\in P$, then $|\Delta u|$
induces homology isomorphisms in every degree.
\end{lemma}

\begin{proof}
Form a poset $M$ on the disjoint union $S\sqcup P$, retaining the two
internal orders and declaring $x\leq p$ for $x\in S$, $p\in P$ precisely
when $u(x)\leq p$. There are no relations from $P$ to $S$. The map
$\rho:M\to P$ equal to $u$ on $S$ and to the identity on $P$ is
order-preserving. For the inclusion $i:P\to M$ we have $\rho i=\mathrm{id}_P$
and $\mathrm{id}_M\leq i\rho$ pointwise.

Comparable order maps induce homotopic realization maps. To see this for
maps $a\leq b$, triangulate the prism on a chain $x_0<\cdots<x_n$ by the
simplices with bottom vertices $x_0,\ldots,x_j$ and top vertices
$x_j,\ldots,x_n$. Their images
$a(x_0),\ldots,a(x_j),b(x_j),\ldots,b(x_n)$ are weak chains; affine extension
on these simplices gives compatible maps on all prism faces. It follows
that $|\Delta M|$ is homotopy equivalent to $|\Delta P|$.

Start with $\Delta S$ and adjoin the elements of $P$ in reverse linear
extension order. When $p$ is added, all elements of $P$ strictly above it
are present and all those strictly below it are absent. Its link is
\begin{equation}\label{eq:link}
 L_p*\Delta(P_{>p}),\qquad
 L_p=\Delta\{x\in S:u(x)\leq p\}.
\end{equation}
Indeed these are exactly its lower and upper elements in the current poset,
and every such lower element is below every such upper element.

The join of a nonempty acyclic finite complex $L$ with any finite complex
$U$ is acyclic. We include the convention that the complex with only the
empty simplex is the identity for join. Using augmented chains, with the
empty simplex in degree $-1$, we have
\[
 \widetilde C_n(L*U;\F)
 =\bigoplus_{\ell+k=n-1}\widetilde C_\ell(L;\F)\otimes
                  \widetilde C_k(U;\F).
\]
The differential is $d_L\otimes1+1\otimes d_U$ after this degree shift.
The exact finite vector-space complex $\widetilde C_*(L;\F)$ admits a
contraction $h$: choose complements to its boundary subspaces and invert
the differential from each complement onto the preceding boundary space.
Then $d_Lh+hd_L=1$. On the join chains, $H=h\otimes1$ satisfies
\[
 dH+Hd=(d_Lh+hd_L)\otimes1+2h\otimes d_U=1,
\]
since the coefficients have characteristic two. This proves the assertion,
including when the upper poset in \eqref{eq:link} is empty.

Adding $p$ attaches the cone on its link. The new relative chains are the
augmented chains of that link shifted up one degree: deleting $p$ gives
a face in the old complex, and the remaining boundary is the link's
augmented boundary. Their homology is zero, so the long exact sequence of
the pair makes each attachment a homology isomorphism
\cite[Theorem~2.16 and pp.~117--118, 153]{Hatcher}. Thus
$\Delta S\hookrightarrow\Delta M$ induces homology isomorphisms.
Composing with $\rho$ proves the claim for $u$.
\end{proof}

\begin{proof}[Proof of Theorem~\ref{thm:main}]
For nonempty $K$, let $S(K)$ be its face poset of nonempty simplices, and
define $u:S(K)\to\A(K)$ by sending $\sigma$ to the subcomplex of $\sigma$
and all its faces. This subcomplex is a simplex, hence acyclic, and $u$ is
order-preserving. For $A\in\A(K)$, the condition $u(\sigma)\subseteq A$
is equivalent to $\sigma\in A$. The lower fiber in Lemma~\ref{lem:poset}
is therefore the face poset of $A$. Its order complex is the barycentric
subdivision of $A$, so it is nonempty and acyclic.

Lemma~\ref{lem:poset} and the barycentric subdivision of $K$ now give
\[
 H_q(|\Delta\A(K)|;\F)\cong H_q(K;\F)
 \qquad(q\geq0).
\]
Proposition~\ref{prop:chains} identifies the space on the left with $G_K$,
and Lemma~\ref{lem:truncate} replaces $G_K$ by $F_K$ without changing its
homology. This proves the asserted isomorphism for every $q\geq0$.

If the convention for connected complexes permits an empty complex, its
unique filter has empty barcode and the prescribed fiber is empty. Ordinary
homology of both empty spaces is zero in every nonnegative degree, so the
same conclusion holds in that case.
\end{proof}

\paragraph{AI use.}
OpenAI Codex was used to develop the argument, check the proof and references,
and prepare this note. Independent mathematical review is requested.

\begin{thebibliography}{9}
\bibitem{BWW}
A.~Bj\"orner, M.~L.~Wachs and V.~Welker,
\emph{Poset fiber theorems},
Transactions of the American Mathematical Society \textbf{357} (2005),
1877--1899. \url{https://doi.org/10.1090/S0002-9947-04-03496-8}.
Author manuscript: \url{https://www.mathematik.uni-marburg.de/~welker/preprints/fiber.pdf}.

\bibitem{HTT}
H.~A.~Harrington, U.~Tillmann and M.~Torras-P\'erez,
\emph{The fiber of multiparameter persistent homology for simplicial complexes},
arXiv:2608.04640v1 (2026). \url{https://arxiv.org/abs/2608.04640v1}.

\bibitem{Hatcher}
A.~Hatcher, \emph{Algebraic Topology},
Cambridge University Press, 2002.

\bibitem{LH}
J.~Leygonie and G.~Henselman-Petrusek,
\emph{Algorithmic reconstruction of the fiber of persistent homology on cell complexes},
Journal of Applied and Computational Topology \textbf{8} (2024), 2015--2049.
\url{https://doi.org/10.1007/s41468-024-00165-w}.

\bibitem{Munkres}
J.~R.~Munkres, \emph{Topology}, second edition,
Prentice Hall, 2000.
\end{thebibliography}
\end{document}
